\SourcePage{33}
\SGAsection{3}{对偶化复形的存在性}
\textbf{3.1. 预备知识。}
如引言中所说，本节的目的是证明：在一个“良好”的\Term{regular}{正则}\Term{prescheme}{预概形} $X$ 上，\Term{sheaf}{层} $\Lambda_X$ 是\Term{dualizing-complex}{对偶化复形}；这里假定 $n$ 与各\Term{residue-characteristic}{剩余特征}互素。实际上，我们希望在\Term{equal-characteristic}{等特征}的\Term{excellent}{优秀}\Term{regular}{正则}\Term{prescheme}{预概形}上也有这一结论；毫无疑问，一旦在这种情形中具备\Term{resolution-singularities}{奇点消解}和\Term{purity}{纯性}\Term{theorem}{定理}，下述证明就会推广到这些\Term{prescheme}{预概形}。
我们先介绍将要用到的各项工具。

\textbf{3.1.1. \Term{singular-locus}{奇异点集}。} 设 $X$ 为（\Term{locally-noetherian}{局部诺特}的）\Term{prescheme}{预概形}。若下列\Term{equivalence}{等价}性质成立（EGA IV 6.12.3），就称 $X$ 满足条件 $(\mathrm{reg})$：
\begin{enumerate}
\item[(i)] 对 $X$ 的每个子\Term{prescheme}{预概形} $Y$，$Y$ 的\Term{regular}{正则}点集 $\operatorname{Reg}(Y)$ 在 $Y$ 中是开的；
\item[(ii)] 对 $X$ 的每个\Term{integral}{整}闭子\Term{prescheme}{预概形} $Y$，$\operatorname{Reg}(Y)$ 包含 $Y$ 的一个\Term{nonempty}{非空}\Term{open-subset}{开集}。
\end{enumerate}
若 $X$ 在一个\Term{complete-noetherian-local-ring}{完备诺特局部环}（例如域）上\Term{locally-finite-type}{局部有限型}（EGA IV 6.12.8），或在\Term{field-of-fractions}{分式域}\Term{characteristic}{特征}为零的 Dedekind 环（例如 $\mathbb Z$）上\Term{locally-finite-type}{局部有限型}（EGA IV 6.12.6），则条件 $(\mathrm{reg})$ 成立；更一般地，若 $X$ 是\Term{excellent}{优秀}的，这一条件也成立（EGA IV 7.8.6）。

\SourcePage{34}
\textbf{3.1.2. \Term{constructible-sheaf}{可构造层}的\Term{resolution}{分解}。} 设 $X$ 为\Term{noetherian}{诺特}\Term{prescheme}{预概形}，$C$ 为\Term{constructible}{可构造} $\Lambda$-\Term{module-category}{模范畴} $\operatorname{Cons}(X)$ 的\Term{strictly-full-subcategory}{严格满子范畴}。假定 $C$ 对\Term{direct-summand}{直和因子}及\Term{group-extension}{扩张}封闭。以下条件\Term{equivalence}{等价}：
\begin{enumerate}
\item[(i)] $C=\operatorname{Cons}(X)$；
\item[(ii)] $C$ 包含形如 $i_!(F)$ 的\Term{sheaf}{层}，其中 $i:Y\to X$ 是一个\Term{integral}{整}子\Term{prescheme}{预概形}的\Term{immersion}{浸入}，$F$ 在 $Y$ 上\Term{locally-constant}{局部常值}；
\item[(iii)] $C$ 包含形如 $f_*i_!(\Lambda'_U)$ 的\Term{sheaf}{层}，其中 $i:U\to Y$ 是\Term{integral}{整}\Term{prescheme}{预概形}的\Term{open-subset}{开集}，$f:Y\to X$ 是\Term{finite}{有限}\Term{morphism}{态射}，$Y$ 的\Term{generic-point}{泛点}在其像上\Term{separable}{可分}，并且 $\Lambda'=\mathbb Z/\ell\mathbb Z$，$\ell$ 为 $n$ 的素因子。
\end{enumerate}
若 $X$ 满足 $(\mathrm{reg})$（3.1.1），则这些条件还\Term{equivalence}{等价}于：
\begin{enumerate}
\item[(ii bis)] 与 (ii) 相同，但另要求 $Y$ 为\Term{regular}{正则}的；
\item[(iii bis)] 与 (iii) 相同，但要求 $U$ 为\Term{regular}{正则}的。
\end{enumerate}
若 $X$ 还是\Term{affine}{仿射}的，则在 (ii)（分别地 (ii bis)）中可假定 $Y$ 是某个形如 $D(f)$ 的\Term{open-subset}{开集}中的\Term{closed-subset}{闭集}，其中 $f\in\Gamma(X;\mathscr O_X)$。
\begin{proof}
(i)—(iii) 的\Term{equivalence}{等价}性已在 SGA 4 IX 2.5 和 5.8 中证明。其余断言由此通过\Term{noetherian-induction}{诺特归纳}法得到。
\end{proof}
\textbf{3.1.3. \Term{cohomological-dimension}{上同调维数}。} 设 $X$ 为（\Term{locally-noetherian}{局部诺特}的）\Term{prescheme}{预概形}。以 $(\mathrm{cdloc})_n$ 表示以下条件：对 $X$ 的每个子\Term{prescheme}{预概形} $Y$、$Y$ 的每个\Term{geometric-point}{几何点} $\bar y$，以及每个
\SourcePage{35}
$f\in\Gamma(\overline Y(\bar y);\mathscr O_{Y,\bar y})$（这里 $\overline Y(\bar y)$ 表示 $Y$ 在 $\bar y$ 处的\Term{strict-localization}{严格局部化}，参见 SGA 4 VIII 7.1），有
\[
\operatorname{cd}_{L}(\overline Y(\bar y)-V(f))\leq\dim\mathscr O_{Y,\bar y}
\quad(=\dim\mathscr O_{Y,y}\text{，由 EGA IV 6.1.3}),
\]
其中 $L$ 是整除 $n$ 的素数\Term{set}{集合}。
若 $X$ 在域上\Term{locally-finite-type}{局部有限型}（SGA 4 XIV 3.5），或 $X$ 是\Term{characteristic}{特征}为零的\Term{excellent}{优秀}\Term{prescheme}{预概形}（SGA 4 XIX 6.2），则 $(\mathrm{cdloc})_n$ 成立。\footnote{（1977 年增补）或者，由 G. Ofer 的一个近期结果，当 $X$ 为\Term{excellent}{优秀}、\Term{regular}{正则}且\Term{dimension}{维数}为 $2$ 时。}

\textbf{3.1.4. \Term{absolute-cohomological-purity}{绝对上同调纯性}。} 设 $X$ 为\Term{regular}{正则}\Term{prescheme}{预概形}，$Y$ 为 $X$ 的\Term{regular}{正则}闭子\Term{prescheme}{预概形}，处处\Term{codimension}{余维数}为 $d$。若
\[
\mathscr H^i_Y(\Lambda_X)=0\quad(i\ne2d),
\]
且“\Term{local-fundamental-class}{局部基本类}”\Term{homomorphism}{同态}（SGA $4\frac12$ Cycle 2.2）
\[
\Lambda_Y\longrightarrow\mathscr H^{2d}_Y((\mu_n)_X^{\otimes d})
\]
是\Term{isomorphism}{同构}，就称 $(X,Y)$ 满足（相对于 $n$ 的）\Term{absolute-purity}{绝对上同调纯性}\Term{theorem}{定理}。
若对 $X$ 的每个\Term{open-subset}{开集} $U$，上述每一对 $(U,Y)$ 都满足此\Term{theorem}{定理}，就称\Term{absolute-purity}{绝对（上同调）纯性}\Term{theorem}{定理}在 $X$ 上成立。
若 $X$ 在\Term{characteristic}{特征}与 $n$ 互素的\Term{perfect-field}{完美域}上\Term{smooth}{光滑}（SGA 4 XVI 3.9），或为\Term{characteristic-zero}{特征零}的\Term{excellent}{优秀}\Term{prescheme}{预概形}（SGA 4 XIX 3.2），或为\Term{dimension}{维数} $1$ 的\Term{regular}{正则}\Term{prescheme}{预概形}（参见 5.1），则\Term{absolute-purity}{绝对纯性}\Term{theorem}{定理}在 $X$ 上成立。\footnote{（1977 年增补）或者，由 G. Ofer 的一个近期结果，当 $X$ 为\Term{excellent}{优秀}、\Term{regular}{正则}且\Term{dimension}{维数}为 $2$ 时。}

\SourcePage{36}
\textbf{3.1.5. 广中式\Term{resolution-singularities}{奇点消解}。}
a) 设 $X$ 为\Term{regular}{正则}\Term{prescheme}{预概形}，$D$ 为 $X$ 上的一个 $\geq0$ 的\Term{divisor}{除子}。若存在 $\Gamma(X;\mathscr O_X)$ 中元素的\Term{finite}{有限}族 $(f_i)_{i\in I}$，使得
\[
D=\sum_{i\in I}\operatorname{div}(f_i)
\]
（EGA IV 21），并且对每个 $x\in\operatorname{Supp}(D)$，那些满足 $(f_i)_x\in\mathfrak m_x$ 的 $(f_i)_x$ 构成一个\Term{regular-parameters}{正则参数系}的一部分，就称 $D$ 为\Term{strict-normal-crossings}{严格正规交叉}的。如果 $D$ 在\Term{etale}{平展}\Term{topology}{拓扑}下\Term{local-scheme}{局部}为\Term{strict-normal-crossings}{严格正规交叉}，就称 $D$ 为\Term{normal-crossings}{正规交叉}的。
标准\Term{affine}{仿射}\Term{space}{空间}中坐标\Term{hyperplane}{超平面}的一个\Term{locally-finite}{局部有限}并给出\Term{normal-crossings}{正规交叉}\Term{divisor}{除子}的典型例子。

b) 设 $X$ 为\Term{locally-noetherian}{局部诺特}\Term{prescheme}{预概形}。若下列条件成立，就称 $X$ 为\Term{strongly-desingularizable}{强可消解}的：对任意\Term{finite}{有限}\Term{morphism}{态射} $f:Y\to X$，其中 $Y$ 为\Term{integral}{整}的且其\Term{generic-point}{泛点}在像上\Term{separable}{可分}，以及 $Y$ 的任意\Term{regular}{正则}\Term{open-subset}{开集} $U$，存在\Term{proper}{固有}\Term{birational}{双有理}\Term{morphism}{态射} $h:Y'\to Y$，使 $Y'$ 为\Term{regular}{正则}的，$h$ 诱导 $U'=h^{-1}(U)$ 到 $U$ 的\Term{isomorphism}{同构}，且 $Y'-U'$ 是\Term{normal-crossings}{正规交叉}\Term{divisor}{除子}。
下列情形中的 $X$ 是\Term{strongly-desingularizable}{强可消解}的：
\begin{enumerate}
\item[a)] $X$ 为\Term{characteristic-zero}{特征零}的\Term{excellent}{优秀}\Term{prescheme}{预概形}（[2]）；
\item[b)] $X$ 为\Term{locally-noetherian}{局部诺特}、\Term{regular}{正则}且\Term{dimension}{维数} $\leq1$；
\item[c)] $X$ 在 $\mathbb Z$ 或\Term{perfect-field}{完美域}上\Term{finite-type}{有限型}，且\Term{dimension}{维数} $\leq2$（[1]）。\footnote{（1977 年增补）更一般地，若 $X$ 为\Term{dimension}{维数} $\leq2$ 的\Term{excellent}{优秀}\Term{noetherian}{诺特}\Term{prescheme}{预概形}；参见 H. Hironaka，\emph{\Term{excellent}{优秀}\Term{surface}{曲面}的\Term{resolution-singularities}{奇点消解}}，B. Bennett 记录，Bowdoin College，1967（必要时补以 H. Hironaka，\emph{\Term{singular-point}{奇点}的某些数值\Term{characteristic}{特征}}，J. of Math. of Kyoto Univ.，10，1970，第 151—187 页）。}
\end{enumerate}

\SourcePage{37}
\textbf{3.1.6. \Term{cohomology}{上同调}\Term{finiteness}{有限性}。} 设 $f:X\to Y$ 为\Term{prescheme}{预概形}的\Term{morphism}{态射}，$F\in\operatorname{ob}D_c^b(X)$。在下列情形中有 $\mathrm Rf_*(F)\in\operatorname{ob}D_c^b(Y)$：
\begin{enumerate}
\item[a)] $f$ 为\Term{proper}{固有}且\Term{finite-presentation}{有限呈示}（SGA 4 XIV 1.1）；
\item[b)] “具备\Term{resolution-singularities}{奇点消解}和\Term{purity}{纯性}\Term{theorem}{定理}时”，特别是：
\end{enumerate}
(i) $X$ 和 $Y$ 在\Term{characteristic-zero}{特征零}的域 $k$ 上\Term{locally-finite-type}{局部有限型}，且 $f$ 为\Term{finite-type}{有限型} $k$-\Term{morphism}{态射}（SGA 4 XVI 5.1）；
(ii) $Y$ 为\Term{characteristic-zero}{特征零}的\Term{excellent}{优秀}\Term{prescheme}{预概形}，$f$ 为\Term{finite-type}{有限型}（SGA 4 XIX 5.1）；
(iii) $X$ 在域上\Term{locally-finite-type}{局部有限型}且 $\dim(X)\leq2$（SGA 4 XIX 5.1）。

\textbf{3.2. $\Lambda_X$ 的\Term{quasi-injective-dimension}{拟内射维数}。}
\begin{SGAresult}{命题}{3.2.1}
设 $X$ 为\Term{regular}{正则}\Term{prescheme}{预概形}，各\Term{residue-characteristic}{剩余特征}与 $n$ 互素。假定 $X$ 满足 $(\mathrm{reg})$（3.1.1）与 $(\mathrm{cdloc})_n$（3.1.3），且\Term{absolute-purity}{绝对纯性}\Term{theorem}{定理}在 $X$ 上成立（3.1.4）。则对每个\Term{constructible}{可构造} $\Lambda$-模 $F$ 及每个 $x\in X$，有
\[
\mathscr Ext^i(F,\Lambda_X)_x=0\quad(i>2\dim\mathscr O_{X,x}).\tag{*}
\]
\end{SGAresult}
\begin{proof}
问题是\Term{local-scheme}{局部}的，故可假定 $X$ 为\Term{affine}{仿射}的。通过\Term{resolution}{分解}（3.1.2），只需对 $F=i_!(G)$ 验证 (*)，其中 $i:Y\to X$ 是一个\Term{integral}{整}\Term{regular}{正则}子\Term{prescheme}{预概形}的\Term{immersion}{浸入}，$Y$ 闭于某个\Term{open-subset}{开集} $D(f)$（$f\in A(X)$），而 $G$ 是 $Y$ 上的\Term{locally-constant}{局部常值}\Term{constructible}{可构造}\Term{sheaf}{层}。
\SourcePage{38}
要计算的是 $\mathscr Ext^q(F,\Lambda_X)=\mathscr H^q(\mathrm R\mathscr Hom(F,\Lambda_X))$。$i$ 的\Term{projection-formula}{投影公式}（SGA 4 XVIII，原书此处未填细目）给出
\[
\mathrm R\mathscr Hom(i_!G,\Lambda_X)\simeq\mathrm Ri_*\mathrm R\mathscr Hom(G,\mathrm Ri^!\Lambda_X).
\]
由\Term{purity}{纯性}\Term{theorem}{定理}，$\mathrm Ri^!(\Lambda_X)\simeq(\mu_n)_Y^{\otimes(-d)}[-2d]$，其中 $d=\operatorname{codim}(Y,X)$。置 $T=(\mu_n)_Y^{\otimes(-d)}$。因为 $G$ \Term{locally-constant}{局部常值}且\Term{constructible}{可构造}，对 $Y$ 的每个\Term{geometric-point}{几何点} $\bar y$ 及每个 $p$，有
\[
\mathscr Ext^p(G,T)_{\bar y}\simeq\operatorname{Ext}^p(G_{\bar y},T_{\bar y}).
\]
而 $T_{\bar y}$ 为\Term{injective}{内射}的，因为它\Term{isomorphism}{同构}于 $\Lambda=\mathbb Z/n\mathbb Z$。所以 $p\ne0$ 时 $\mathscr Ext^p(G,T)=0$，从而 $\mathrm R\mathscr Hom(G,T)\simeq\mathscr Hom(G,T)$；最终
\[
\mathrm R\mathscr Hom(F,\Lambda_X)\simeq\mathrm Ri_*(H)[-2d],\qquad H=\mathscr Hom(G,T).
\]
设 $Z$ 是 $Y$ 在 $X$ 中的\Term{schematic-closure}{概形闭包}（EGA I 9.5.10）。有 $Y=Z\cap D(f)$（EGA $0_{\mathrm I}$ 2.1.6）。以 $j:Y\to Z$、$k:Z\to X$ 表示\Term{canonical}{典范}\Term{immersion}{浸入}。有 $\mathrm Ri_*=k_*\mathrm Rj_*$，所以归结为计算 $\mathrm Rj_*(H)$。设 $\bar x$ 为 $Z$ 的\Term{geometric-point}{几何点}，$\overline Z=\operatorname{Spec}(\mathscr O_{Z,\bar x})$ 是 $Z$ 在 $\bar x$ 处的\Term{strict-localization}{严格局部化}，$\bar f$ 是 $f$ 在 $\mathscr O_{Z,\bar x}$ 中的像，$\overline Y=\overline Z\times_XY=\overline Z-V(\bar f)$，$\overline H$ 是 $H$ 在 $\overline Y$ 上的\Term{inverse-image}{逆像}。由 SGA 4 VIII 5.2，
\[
(\mathrm R^qj_*H)_{\bar x}\simeq H^q(\overline Y,\overline H),\qquad
\mathscr Ext^q(F,\Lambda_X)_{\bar x}\simeq H^{q-2d}(\overline Y,\overline H).
\]
条件 $(\mathrm{cdloc})_n$ 因而蕴含 $q-2d>\dim\mathscr O_{Z,x}$ 时 $\mathscr Ext^q(F,\Lambda_X)_{\bar x}=0$。由
\[
d=\operatorname{codim}_x(Z,X)\leq\dim\mathscr O_{X,x}\quad(\text{EGA IV 5.1.3}),
\]
及
\[
\dim\mathscr O_{Z,x}=\dim\mathscr O_{X,x}-\operatorname{codim}_x(Z,X)
\quad(\text{EGA }0_{\mathrm{IV}}\ 16.5.12\text{ 和 IV 5.1.9}),
\]
得到 $q>2\dim\mathscr O_{X,x}$ 时 $\mathscr Ext^q(F,\Lambda_X)_{\bar x}=0$，证毕。
\end{proof}
\textbf{注记 3.2.2。} (3.2.1) 所得的上界是最优的。事实上，设 $x$ 为 $X$ 的一个点，$Y$ 为以 $x$ 为\Term{generic-point}{泛点}的\Term{integral}{整}闭子\Term{prescheme}{预概形}。由于 $X$ 满足 $(\mathrm{reg})$，存在 $Y$ 中包含 $x$ 的\Term{regular}{正则}\Term{open-subset}{开集} $U$。由\Term{purity}{纯性}\Term{theorem}{定理}，
\SourcePage{39}
有 $\mathscr Ext^{2d}(\Lambda_U,\Lambda_X)|U=(\mu_n)_U^{\otimes d}$，其中 $d=\dim\mathscr O_{X,x}=\operatorname{codim}(Y,X)$。
\begin{SGAresult}{推论}{3.2.3}
设 $X$ 为\Term{dimension}{维数} $d$ 的\Term{regular}{正则}\Term{prescheme}{预概形}。若 $X$ 在\Term{characteristic}{特征}与 $n$ 互素的域 $k$ 上\Term{smooth}{光滑}，或为\Term{characteristic-zero}{特征零}的\Term{excellent}{优秀}\Term{prescheme}{预概形}，则 $X$ 满足 (3.2.1) 的结论，因而
\[\operatorname{dim.q.inj}(\Lambda_X)\leq2d.\]
\end{SGAresult}
\begin{proof}
当 $X$ 为\Term{characteristic-zero}{特征零}的\Term{excellent}{优秀}\Term{prescheme}{预概形}，或 $k$ 是\Term{characteristic}{特征}与 $n$ 互素的\Term{perfect-field}{完美域}且 $X$ 在 $k$ 上\Term{smooth}{光滑}时，由 (3.1.1)、(3.1.3) 和 (3.1.4)，(3.2.1) 的假设成立。$X$ 在任意域 $k$ 上\Term{smooth}{光滑}时结论仍然成立：以 $k'$ 为 $k$ 的\Term{perfect-closure}{完美闭包}，作\Term{base-change}{基变换} $\operatorname{Spec}(k')\to\operatorname{Spec}(k)$ 即可看出（SGA 4 VIII 1.1）。
\end{proof}
\begin{SGAresult}{\Term{lemma}{引理}}{3.2.4}
设 $X$ 满足 (3.2.1) 的结论且\Term{dimension}{维数} $\leq d$。若存在整数 $N$ 使 $X$ 的\Term{cohomological-dimension}{上同调维数} $\leq N$，则
\[\operatorname{dim.inj}(\Lambda_X)\leq2d+N.\]
若对 $X$ 的每个\Term{closed-subset}{闭集} $Y$，$\operatorname{codim}(Y,X)\geq p$ 还蕴含 $\operatorname{cd}_n(Y)\leq N-2p$，则
\[\operatorname{dim.inj}(\Lambda_X)\leq N.\]
\end{SGAresult}
\begin{proof}
第一断言立即由 (1.5) 得到。第二断言的论证与 (1.5) 相同。归结为证明：对 $X$ 上\Term{constructible}{可构造}的 $F$，$i>N$ 时 $\operatorname{Ext}^i(X;F,\Lambda_X)=0$。考察\Term{spectral-sequence}{谱序列}
\SourcePage{40}
\[
E_2^{pq}=H^p(X;\mathscr Ext^q(F,\Lambda_X))\Longrightarrow\operatorname{Ext}^{*}(X;F,\Lambda_X).
\]
由 (3.2.1)，\Term{sheaf}{层} $\mathscr Ext^{2i}(F,\Lambda_X)$ 和 $\mathscr Ext^{2i-1}(F,\Lambda_X)$ 集中在\Term{codimension}{余维数} $\geq i$ 处。因而由假设，当 $p>N-2i$ 时 $E_2^{p,2i}=E_2^{p,2i-1}=0$。所以 $p+q>N$ 时 $E_2^{pq}=0$，断言得证。
\end{proof}
\begin{SGAresult}{推论}{3.2.5}
若 $X$ 是在\Term{characteristic}{特征}与 $n$ 互素的\Term{separably-closed-field}{可分闭域} $k$ 上\Term{smooth}{光滑}的 $d$ 维\Term{prescheme}{预概形}，则 $\operatorname{dim.inj}(\Lambda_X)=2d$。
\end{SGAresult}
\begin{proof}
由 (3.2.3) 和 SGA 4 X 4.2，我们确实处于可应用 (3.2.4) 的条件下（注意 SGA 4 X 4.2 的假设由 SGA 4 X 2.1 和 3.2 保证），故 $\operatorname{dim.inj}(\Lambda_X)\leq2d$。等式由\Term{purity}{纯性}\Term{theorem}{定理}得到（参见 3.2.2）：若 $x$ 为 $X$ 中\Term{codimension}{余维数}为 $d$ 的\Term{closed-point}{闭点}，则
\[
\operatorname{Ext}^{2d}(X,\Lambda_x,\Lambda_X)=H_x^{2d}(\Lambda_X)\simeq\mu_n^{\otimes d}(k)\ne0.
\]
\end{proof}

\textbf{3.3. $\mathscr Ext^i(F,\Lambda_X)$ 的\Term{constructibility}{可构造性}。}
\begin{SGAresult}{命题}{3.3.1}
设 $X$ 为\Term{prescheme}{预概形}。
a) 下列条件\Term{equivalence}{等价}：
\begin{enumerate}
\item[(i)] 对 $F\in\operatorname{ob}D_c^-(X)$、$G\in\operatorname{ob}D_c^+(X)$，$\mathrm R\mathscr Hom(F,G)\in\operatorname{ob}D_c^+(X)$；
\item[(i bis)] 对任意一对\Term{constructible}{可构造} $\Lambda_X$-模 $(F,G)$，$\mathrm R\mathscr Hom(F,G)\in\operatorname{ob}D_c^+(X)$，即对每个 $i$，$\mathscr Ext^i(F,G)$ 为\Term{constructible}{可构造}的；
\item[(ii)] 对每个\Term{open-immersion}{开浸入} $i:U\to X$ 及每个\Term{constructible}{可构造} $\Lambda_U$-模 $F$，$\mathrm R^pi_*(F)$ 对每个 $p$ 均为\Term{constructible}{可构造}的；
\SourcePage{41}
\item[(ii bis)] 对每个\Term{immersion}{浸入} $i:Y\to X$ 及每个\Term{constructible}{可构造} $\Lambda_Y$-模 $F$，$\mathrm R^pi_*(F)$ 对每个 $p$ 均为\Term{constructible}{可构造}的；
\item[(ii ter)] 对每个\Term{immersion}{浸入} $i:Y\to X$ 及每个\Term{locally-constant}{局部常值}\Term{constructible}{可构造} $\Lambda_Y$-模 $F$，$\mathrm R^pi_*(F)$ 对每个 $p$ 均为\Term{constructible}{可构造}的；
\item[(iii)] 对每个\Term{finite}{有限}\Term{morphism}{态射}（分别地，每个\Term{closed-immersion}{闭浸入}）$f:Y\to X$，以及每个 $F\in\operatorname{ob}D_c^+(X)$，有 $\mathrm Rf^!(F)\in\operatorname{ob}D_c^+(Y)$。
\end{enumerate}
b) 在以下两种情形中这些条件成立：
(i) $\dim(X)\leq1$；
(ii) “在 $X$ 上具备\Term{resolution-singularities}{奇点消解}与\Term{purity}{纯性}”（参见 3.1.6 b)）。
\end{SGAresult}
\begin{proof}
(i) 与 (i bis) 的\Term{equivalence}{等价}性来自从 Ext 到超 Ext 的标准\Term{spectral-sequence}{谱序列}。
(i bis)$\Rightarrow$(ii)。设 $\overline F$ 为延拓 $F$ 的\Term{constructible}{可构造} $\Lambda_X$-模，例如 $\overline F=i_!(F)$。由\Term{duality}{对偶性}（SGA 4 XVIII 3.1.10），有
\[
\mathrm Ri_*(F)=\mathrm Ri_*\mathrm R\mathscr Hom(\Lambda_U,i^!\overline F)
\simeq\mathrm R\mathscr Hom(i_!\Lambda_U,\overline F),
\]
结论随之得到。
(ii)$\Rightarrow$(ii bis)。把 $i$ \Term{resolution}{分解}为一个\Term{closed-immersion}{闭浸入}之后接一个\Term{open-immersion}{开浸入}。
(ii bis)$\Rightarrow$(i bis)。问题是\Term{local-scheme}{局部}的，故可假定 $X$ 为\Term{affine}{仿射}的。由通常的\Term{resolution}{分解}（SGA 4 IX 2.5），只需对 $F=i_!(M)$ 验证 (ii)；其中 $i:Y\to X$ 是子\Term{prescheme}{预概形}的\Term{immersion}{浸入}，$M$ 为\Term{locally-constant}{局部常值}（\Term{constructible}{可构造}）$\Lambda_Y$-模。由\Term{projection-formula}{投影公式}
\SourcePage{42}
（SGA 4 XVIII，原书此处未填细目）
\[
\mathrm R\mathscr Hom(i_!M,G)\simeq\mathrm Ri_*\mathrm R\mathscr Hom(M,\mathrm Ri^!G)
\]
及一个\Term{spectral-sequence}{谱序列}，归结为证明 $\mathrm R\mathscr Hom(M,\mathrm Ri^!G)\in\operatorname{ob}D_c^+(Y)$。
先证明 $\mathrm Ri^!G\in\operatorname{ob}D_c^+(Y)$。将 $i$ \Term{resolution}{分解}后，可以假定 $Y$ 为闭的。设 $j:U=X-Y\to X$ 为互补\Term{open-immersion}{开浸入}。对\Term{exact-sequence}{正合列}
\[0\longrightarrow\Lambda_{U,X}\longrightarrow\Lambda_X\longrightarrow\Lambda_Y\longrightarrow0\]
施以\Term{functor}{函子} $\mathrm R\mathscr Hom(-,G)$，得到\Term{distinguished-triangle}{特殊三角形}
\[
\begin{tikzcd}
&\mathrm R\mathscr Hom(\Lambda_{U,X},G)\arrow[dl]\\
\mathrm R\mathscr Hom(\Lambda_Y,G)\arrow[rr]&&\mathrm R\mathscr Hom(\Lambda_X,G)\arrow[ul]
\end{tikzcd}
\]
而 $\mathrm R\mathscr Hom(\Lambda_X,G)=G$；另一方面，由假设，$\mathrm R\mathscr Hom(\Lambda_{U,X},G)\simeq\mathrm Rj_*j^!G$ 属于 $D_c^+(X)$。所以 $\mathrm R\mathscr Hom(\Lambda_Y,G)\simeq i_*\mathrm Ri^!G$ 属于 $D_c^+(X)$。在\Term{etale}{平展}\Term{topology}{拓扑}下\Term{local-scheme}{局部}用自由 $\Lambda_Y$-模左\Term{resolution}{分解} $M$，再作一次\Term{resolution}{分解}便得到 $\mathrm R\mathscr Hom(M,\mathrm Ri^!G)\in\operatorname{ob}D_c^+(X)$，证毕。
(ii bis)$\Rightarrow$(ii ter) 显然。
(ii ter)$\Rightarrow$(ii)。设 $i:U\to X$ 为\Term{open-immersion}{开浸入}，$F$ 为\Term{constructible}{可构造} $\Lambda_U$-模。须证明 $\mathrm Ri_*F\in\operatorname{ob}D_c^+(X)$。问题是\Term{local-scheme}{局部}的，故可假定 $X$ 为\Term{noetherian}{诺特}的，并对 $X$ 作\Term{noetherian-induction}{诺特归纳}。可假定 $U\ne\varnothing$，因而存在
\SourcePage{43}
$U$ 的\Term{nonempty}{非空}\Term{open-subset}{开集} $j:V\to U$，使 $j^*F$ 为\Term{locally-constant}{局部常值}（SGA 4 IX 2.4）。设 $M$ 为 $F\to\mathrm Rj_*j^*F$ 的\Term{mapping-cylinder}{映射柱}。施以 $\mathrm Ri_*$，得到\Term{distinguished-triangle}{特殊三角形}
\[
\begin{tikzcd}
&\mathrm Ri_*M\arrow[dl,"{+1}"']\\
\mathrm Ri_*F\arrow[rr]&&\mathrm Ri_*\mathrm Rj_*j^*F\arrow[ul]
\end{tikzcd}
\]
由条件 (ii ter)，
\[
\mathrm Ri_*\mathrm Rj_*j^*F\simeq\mathrm R(ij)_*j^*F\in\operatorname{ob}D_c^+(X).
\]
从而 $\mathrm Rj_*j^*F\simeq i^*\mathrm R(ij)_*j^*F\in\operatorname{ob}D_c^+(U)$。由于 $F\in\operatorname{ob}D_c^+(U)$，有 $M\in\operatorname{ob}D_c^+(U)$。此外，$M$ 的\Term{cohomology}{上同调}\Term{support}{支撑}在 $U$ 的一个真\Term{closed-subset}{闭集} $Z$ 中。对 $U\cap\overline Z\to\overline Z$ 应用归纳假设，得到 $\mathrm Ri_*M\in\operatorname{ob}D_c^+(X)$，故 $\mathrm Ri_*F\in\operatorname{ob}D_c^+(X)$。
条件 (i)—(ii ter) 与条件 (iii) 的\Term{equivalence}{等价}性留给读者（利用 1.14）。
b) (i) 通过\Term{normalization}{正规化}，归结到 $X$ 为\Term{regular}{正则}的情形，再利用\Term{purity}{纯性}（将在 5.1 中证明）。
b) (ii) 参见 3.1.6 b)。
\end{proof}
\SourcePage{44}
\textbf{3.4. \Term{regular}{正则}\Term{prescheme}{预概形}上的\Term{biduality-theorem}{双对偶定理}。}
\begin{SGAresult}{\Term{theorem}{定理}}{3.4.1}
（参见 [H] V 2.2。）设 $X$ 是\Term{finite-dimensional}{有限维}的\Term{regular}{正则}（\Term{locally-noetherian}{局部诺特}）\Term{prescheme}{预概形}。假定 $X$ \Term{strongly-desingularizable}{强可消解}（3.1.5），满足 $(\mathrm{reg})$（3.1.1）及 $(\mathrm{cdloc})_n$（3.1.3），而且\Term{absolute-purity}{绝对纯性}\Term{theorem}{定理}（3.1.4）在每个在 $X$ 上\Term{smooth}{光滑}的\Term{prescheme}{预概形}上成立。则 $\Lambda_X$ 是\Term{dualizing-complex}{对偶化复形}（1.6）。
\end{SGAresult}
\begin{proof}
须证明 $\Lambda_X$ 满足 (1.6) 的条件 (i)—(iii)。条件 (i) 由 (3.2.1) 成立；由 (3.3.1 b))，$\Lambda_X$ 也满足 (ii)。剩下条件 (iii)。换言之，由 (1.7)，须证明对每个\Term{constructible}{可构造} $\Lambda_X$-模 $F$，\Term{canonical-homomorphism}{典范同态} $F\to D_XD_XF$ 是\Term{isomorphism}{同构}，其中置 $D_X=\mathrm R\mathscr Hom(-,\Lambda_X)$。问题是\Term{local-scheme}{局部}的，故可假定 $X$ 为\Term{noetherian}{诺特}的；由\Term{resolution}{分解}（3.1.2 (iii bis)），只需考虑形如 $F=f_*i_!(\Lambda'_U)$ 的 $F$，其中 $i:U\to Y$ 是\Term{integral}{整}\Term{prescheme}{预概形}中\Term{regular}{正则}\Term{open-subset}{开集}的包含，$f:Y\to X$ 是\Term{finite}{有限}\Term{morphism}{态射}，使 $Y$ 的\Term{generic-point}{泛点}在其像上\Term{separable}{可分}，$\Lambda'=\mathbb Z/\ell\mathbb Z$，$\ell$ 为 $n$ 的素因子。
由于 $X$ 为\Term{strongly-desingularizable}{强可消解}的，存在（3.1.6）\Term{proper}{固有}\Term{birational}{双有理}\Term{morphism}{态射} $h:Y'\to Y$，使 $Y'$ 为\Term{regular}{正则}的，$h$ 诱导 $U'=h^{-1}(U)$ 到 $U$ 的\Term{isomorphism}{同构}，而 $Y'-U'$ 是\Term{normal-crossings}{正规交叉}\Term{divisor}{除子}（3.1.6）。由\Term{base-change}{基变换}\Term{theorem}{定理}（参见 SGA 4 XVII，原书未填细目），
\[
i_!\Lambda'_U\simeq\mathrm Rh_*i'_!\Lambda'_{U'},
\]
其中 $i'$ 为包含 $U'\to Y'$。故
\[
f_*i_!\Lambda'_U\simeq f_*\mathrm Rh_*i'_!\Lambda'_{U'}\simeq\mathrm Rg_*i'_!\Lambda'_{U'},\quad g=fh.
\]
由 (1.11)，只需证明这一对
\SourcePage{45}
$(i'_!\Lambda'_{U'},\mathrm Rg^!\Lambda_X)$ 为\Term{bidualizing}{双对偶化}的。在\Term{local-scheme}{局部}把 $Y'$ \Term{immersion}{浸入}一个在 $X$ 上\Term{smooth}{光滑}的\Term{prescheme}{预概形} $X'$，并在 $X'$ 上应用\Term{purity}{纯性}\Term{theorem}{定理}，得到 $\mathrm Rg^!\Lambda_X$ \Term{local-isomorphism}{局部同构}于 $\Lambda_{Y'}$ 的某个\Term{translation}{平移}。因此归结为证明 (3.4.1) 的如下特殊情形。
\end{proof}
\begin{SGAresult}{\Term{lemma}{引理}}{3.4.2}
设 $X$ 为\Term{locally-noetherian}{局部诺特}\Term{regular}{正则}\Term{prescheme}{预概形}，\Term{absolute-purity}{绝对纯性}\Term{theorem}{定理}在其上成立（3.1.4）。设 $Y=Y_1\cup\cdots\cup Y_p$ 为 $X$ 上\Term{strict-normal-crossings}{严格正规交叉}\Term{divisor}{除子}（3.1.6），$U=X-Y$。则 $(\Lambda'_{U,X},\Lambda_X)$ 为\Term{bidualizing}{双对偶化}的（1.4）。
\end{SGAresult}
\begin{proof}
由\Term{exact-sequence}{正合列} $0\to\Lambda'_{U,X}\to\Lambda'_X\to\Lambda'_Y\to0$，\Term{equivalence}{等价}于验证 $(\Lambda'_Y,\Lambda_X)$ 为\Term{bidualizing}{双对偶化}的。对 $p$ 归纳。
a) $p=1$。记 $i:Y\to X$ 为\Term{immersion}{浸入}。由 (1.11)，只需验证 $(\Lambda'_Y,\mathrm Ri^!\Lambda_X)$ 为\Term{bidualizing}{双对偶化}的。而由\Term{purity}{纯性}\Term{theorem}{定理}，$\mathrm Ri^!\Lambda_X\simeq(\mu_n)_Y^{\otimes(-1)}[-2]$。从而\Term{canonical-homomorphism}{典范同态} $\Lambda'_Y\to D_YD_Y(\Lambda'_Y)$ 是\Term{isomorphism}{同构}；逐\Term{fiber}{纤维}验证即可（$\Lambda$-模的 Pontryagin \Term{duality}{对偶性}）。
b) $p\geq2$。置 $Y'=Y_1\cup\cdots\cup Y_{p-1}$。有\Term{exact-sequence}{正合列}
\[
0\longrightarrow\Lambda'_{(Y-Y'),X}\longrightarrow\Lambda'_Y\longrightarrow\Lambda'_{Y'}\longrightarrow0.
\]
由归纳假设，$(\Lambda'_{Y'},\Lambda_X)$ 为\Term{bidualizing}{双对偶化}的，故由\Term{five-lemma}{五引理}归结为验证 $(\Lambda'_{(Y-Y'),X},\Lambda_X)$ 的\Term{biduality}{双对偶性}。又有\Term{exact-sequence}{正合列}
\[
0\longrightarrow\Lambda'_{(Y-Y'),X}\longrightarrow\Lambda'_{Y_p}\longrightarrow\Lambda'_{Y_p\cap Y'}\longrightarrow0.
\]
\SourcePage{46}
由于 $(\Lambda'_{Y_p},\Lambda_X)$ 为\Term{bidualizing}{双对偶化}的，只须对 $(\Lambda'_{Y_p\cap Y'},\Lambda_X)$ 验证同样的性质。再应用 (1.11) 和\Term{purity}{纯性}，最终归结为在 $Y_p$ 上验证 $(\Lambda'_{Y_p\cap Y'},\Lambda_{Y_p})$ 的\Term{biduality}{双对偶性}，而这由归纳假设得到。
\end{proof}
\begin{SGAresult}{推论}{3.4.3}
设 $S$ 为\Term{characteristic-zero}{特征零}的\Term{noetherian}{诺特}\Term{regular}{正则}\Term{excellent}{优秀}\Term{prescheme}{预概形}。则 $\Lambda_S$ 是\Term{dualizing-complex}{对偶化复形}。此外，若 $f:X\to S$ 为\Term{finite-type}{有限型}\Term{morphism}{态射}，则 $\mathrm Rf^!(\Lambda_S)$ 是\Term{dualizing-complex}{对偶化复形}。
\end{SGAresult}
\begin{proof}
第一断言来自 $S$ 满足 (3.4.1) 的假设（见 3.1）。第二断言先须验证 $\mathrm Rf^!\Lambda_S$ 的\Term{quasi-injective-dimension}{拟内射维数}\Term{finite}{有限}。为此，用\Term{finite}{有限}个\Term{affine}{仿射}\Term{open-subset}{开集} $S_i$ \Term{covering}{覆盖} $S$，使 $X_i=f^{-1}(S_i)$ 为\Term{finite}{有限}个\Term{affine}{仿射}\Term{open-subset}{开集} $X_{ij}$ 之并，且各 $X_{ij}$ 作为\Term{closed-subset}{闭集}\Term{immersion}{浸入}在 $S_i$ 上\Term{finite-dimensional}{有限维}\Term{smooth}{光滑}的 $X'_{ij}$ 中。由 (3.2.3) 及 (1.6)，
\[
\operatorname{dim.q.inj}(\mathrm Rf^!\Lambda_S)\leq2\sup\dim X'_{ij}.
\]
剩下验证 (1.7) 的 (ii)、(iii)。这些条件在 $X$ 上对\Term{etale}{平展}\Term{topology}{拓扑}是\Term{local-scheme}{局部}的，故可假定 $f=f'i$，其中 $f':X'\to S$ 为\Term{relative-dimension}{相对维数} $d$ 的\Term{smooth}{光滑}\Term{morphism}{态射}，$i:X\to X'$ 为\Term{closed-immersion}{闭浸入}。有 $\mathrm Rf^!\Lambda_S\simeq\mathrm Ri^!\mathrm Rf'^!\Lambda_S$。而 $\mathrm Rf'^!\Lambda_S$ \Term{local-isomorphism}{局部同构}于 $\Lambda_{X'}$（相差 $2d$ 的\Term{translation}{平移}），故由第一断言为\Term{dualizing-complex}{对偶化复形}。由 (1.15)，$\mathrm Rf^!\Lambda_S$ 为\Term{dualizing-complex}{对偶化复形}，特别满足 (1.7) 的条件 (ii) 和 (ii)，证毕。
\end{proof}
\begin{SGAresult}{推论}{3.4.4}
设 $k$ 为域，$f:X\to S=\operatorname{Spec}(k)$ 为\Term{locally-finite-type}{局部有限型}\Term{morphism}{态射}，且 $\dim(X)\leq2$。
\SourcePage{47}
则 $\mathrm Rf^!\Lambda_S$ 为\Term{dualizing-complex}{对偶化复形}。
\end{SGAresult}
\begin{proof}
把 $k$ 换为其\Term{perfect-closure}{完美闭包}，这一操作不影响\Term{etale}{平展}\Term{topology}{拓扑}（VIII 1.1），故可假定 $k$ \Term{perfect}{完美}。用在 $k$ 上\Term{finite-type}{有限型}的\Term{affine}{仿射}\Term{open-subset}{开集} $X_i$ \Term{covering}{覆盖} $X$。因 $\dim X_i\leq2$，由\Term{normalization}{正规化}\Term{lemma}{引理}，$X_i$ 在某个在 $S$ 上\Term{smooth}{光滑}且\Term{dimension}{维数} $\leq2$ 的 $X'_i$ 上\Term{finite}{有限}。由 (1.15)，归结到 $f$ \Term{smooth}{光滑}的情形；此时\Term{local-scheme}{局部}有 $f^!\Lambda_S\simeq\Lambda_X[2d]$，即须证明此时 $\Lambda_X$ 为\Term{dualizing-complex}{对偶化复形}。这由 (3.4.1) 得到，因为各项假设均成立（见 3.1）。
\end{proof}
\textbf{注记 3.4.5。} 下文将看到，若 $X$ 为\Term{dimension}{维数} $\leq1$ 的\Term{regular}{正则}\Term{prescheme}{预概形}，则 $\Lambda_X$ 是\Term{dualizing-complex}{对偶化复形}。

\SGAsection{4}{局部对偶性}
本节的目的是以一种更接近几何直观的\Term{form}{形式}表述前几节发展的“\Term{biduality}{双对偶性}”概念。
\textbf{4.1.} 设 $X$ 为\Term{prescheme}{预概形}，$i:x\to X$ 为\Term{residue-field}{剩余域}\Term{separably-closed}{可分闭}的\Term{closed-point}{闭点}，$K_X$ 为 $X$ 上的\Term{dualizing-complex}{对偶化复形}（1.7）。由 (1.15)，$K_{\{x\}}=\mathrm Ri^!K_X$ 是 $x$ 上的\Term{dualizing-complex}{对偶化复形}；由 (2.1) 及 (3.4.4)，它\Term{isomorphism}{同构}于 $\Lambda[r]$，其中 $r\in\mathbb Z$（若 $\Lambda$ 为\Term{local-ring}{局部环}）。若 $r=0$ 且给定了\Term{isomorphism}{同构} $K_{\{x\}}\simeq\Lambda$，就称 $K_X$ 在 $x$ 处\Term{normalized}{正规化}。
\SourcePage{48}
例如，若 $X$ 在一个\Term{separably-closed-field}{可分闭域}的谱上\Term{smooth}{光滑}，\Term{relative-dimension}{相对维数}为 $d$，则由\Term{purity}{纯性}\Term{theorem}{定理}（SGA XVI 3，SGA $4\frac12$ Cycle 2.3），\Term{complex}{复形} $(\mu_n)_X^{\otimes d}[2d]$ 在 $X$ 的每个\Term{closed-point}{闭点}处\Term{canonical}{典范}地\Term{normalized}{正规化}。

\textbf{4.2.} 设 $X$ 为\Term{prescheme}{预概形}，$i:Y\to X$ 为闭子\Term{prescheme}{预概形}，$K_X$ 为 $X$ 上的\Term{dualizing-complex}{对偶化复形}。则由 (1.15)，$K_Y=\mathrm Ri^!K_X$ 为 $Y$ 上的\Term{dualizing-complex}{对偶化复形}，并有互补\Term{induction-formula}{诱导公式}（1.12 b) (ii)）：对 $F\in\operatorname{ob}D_c^b(X)$，
\[
i^*D_X(F)\simeq D_Y\mathrm Ri^!(F),\tag{4.2.1}
\]
其中置 $D_X=\mathrm R\mathscr Hom(-,K_X)$、$D_Y=\mathrm R\mathscr Hom(-,K_Y)$。
特别把 (4.2.1) 应用于 $Y=\{x\}$ 是 $X$ 中\Term{residue-field}{剩余域}\Term{separably-closed}{可分闭}的\Term{closed-point}{闭点}、且 $K_X$ 在 $x$ 处\Term{normalized}{正规化}（4.1）的情形。由于 $\Lambda_x$ 为\Term{injective}{内射}的，(4.2.1) 写成
\[
D_X(F)_x\simeq\operatorname{Hom}^{\bullet}(\mathrm Ri^!F,\Lambda).\tag{4.2.2}
\]
注意由定义 $\mathrm Ri^!F$ 也写成 $\mathrm R\Gamma_x(F)$，因此在 $D(\Lambda)$ 中有\Term{perfect-pairing}{完美配对}
\[
D_X(F)_x\times\mathrm R\Gamma_x(F)\longrightarrow\Lambda,\tag{4.2.2$'$}
\]
并诱导\Term{finite-group}{有限群}之间的\Term{perfect-pairing}{完美配对}
\[
\mathscr H^i(D_X(F))_x\times H_x^{-i}(F)\longrightarrow\Lambda.\tag{4.2.2$''$}
\]
例如，设 $X$ 为在\Term{separably-closed-field}{可分闭域}上\Term{smooth}{光滑}的 $d$ 维\Term{prescheme}{预概形}，$x\in X$ 为\Term{closed-point}{闭点}。对 $F\in\operatorname{ob}D_c^b(X)$，有\Term{finite-group}{有限群}之间的\Term{canonical}{典范}\Term{perfect-pairing}{完美配对}
\[
\mathscr Ext^{2d-i}(F,(\mu_n)_X^{\otimes d})_x\times H_x^i(F)\longrightarrow\Lambda.\tag{4.2.2$'''$}
\]
\SourcePage{49}
\textbf{习题 4.2.3。} 设 $X$ 为\Term{strictly-local}{严格局部}\Term{scheme}{概形}（SGA 4 VIII 4.2），\Term{closed-point}{闭点}为 $x$，$i:x\to X$ 为\Term{canonical}{典范}\Term{morphism}{态射}。设 $K$ 为 $X$ 上在 $x$ 处\Term{normalized}{正规化}的\Term{dualizing-complex}{对偶化复形}。记 $\mathrm Li^?:D^b(\Lambda)\to D^-(X)$ 为由
\[\mathrm Li^?(M)=M\mathbin{\mathop\otimes\limits^{\mathrm L}_{\Lambda}}K\]
定义的\Term{functor}{函子}。（若 $K$ 在 $\Lambda$ 上 \Term{tor-dimension}{Tor 维数}\Term{finite}{有限}，则 $\mathrm Li^?$ 延拓为从 $D^+(\Lambda)$ 到 $D(X)$ 的\Term{functor}{函子}。）证明对 $L\in D^-(X)$、$M\in D^b(\Lambda)$，存在\Term{canonical}{典范}的\Term{functorial}{函子性}\Term{isomorphism}{同构}
\[
\mathrm R\operatorname{Hom}(L,\mathrm Li^?M)\xrightarrow\sim
\mathrm R\operatorname{Hom}(\mathrm Ri^!L,M).\tag{*}
\]
（当然假定 $\mathrm Ri^!:D^-(X)\to D^-(\Lambda)$ 有定义……。）
\emph{提示：}先定义“迹”\Term{morphism}{态射} $\mathrm Ri^!\mathrm Li^?(M)\to M$。再为证明 (*) 为\Term{isomorphism}{同构}，通过\Term{resolution}{分解}归结到 $L\in\operatorname{ob}D_c^b(X)$、$M=\Lambda$ 的情形，应用 (4.2.2)。注意当 $M=\Lambda$ 时，(*) 化为\Term{perfect-pairing}{完美配对}（“\Term{local-duality}{局部对偶性}”，参见 [H] V 6.2）：
\[
\mathrm R\Gamma_x(L)\times\mathrm R\operatorname{Hom}(L,K)\longrightarrow\Lambda,
\]
或
\[H_x^i(L)\times\operatorname{Ext}^{-i}(L,K)\longrightarrow\Lambda.\]

\textbf{4.3.} 现在将 (4.2.2) 推广到位于 $X$ 的一个不必为\Term{closed-point}{闭点}的点 $x$ 上的\Term{geometric-point}{几何点} $\bar x$。为此需要一些预备结果。
\SourcePage{50}
\begin{SGAresult}{\Term{lemma}{引理}}{4.3}
设 $X$ 为\Term{prescheme}{预概形}，$X'$ 为 $X$ 在\Term{geometric-point}{几何点} $\bar x$ 处的\Term{strict-localization}{严格局部化}，$f:X'\to X$ 为\Term{canonical}{典范}\Term{arrow}{箭头}。则对 $F\in\operatorname{ob}D_c^-(X)$、$G\in\operatorname{ob}D^+(X)$，\Term{canonical}{典范}的\Term{functorial}{函子性}\Term{homomorphism}{同态}（SGA 4 VI 3.7）
\[
f^*\mathrm R\mathscr Hom(F,G)\longrightarrow\mathrm R\mathscr Hom(f^*F,f^*G)
\]
是\Term{isomorphism}{同构}。
\end{SGAresult}
\begin{proof}
照常归结到 $F$ 和 $G$ 都集中在\Term{degree}{次数}零的情形。问题在 $x$ \Term{neighborhood}{邻域}内是\Term{local-scheme}{局部}的，故可假定 $X$ 为\Term{affine}{仿射}的。此时 $F$ 有由形如 $i_!\Lambda_U$ 的\Term{sheaf}{层}组成的左\Term{resolution}{分解}（SGA 4 IX 2.7），其中 $i:U\to X$ 为\Term{finite-presentation}{有限呈示}\Term{etale}{平展}\Term{morphism}{态射}；故通过\Term{resolution}{分解}归结到 $F=i_!\Lambda_U$。作\Term{cartesian-square}{笛卡尔方块}
\[
\begin{tikzcd}
U\arrow[d,"i"']&U'\arrow[l,"g"']\arrow[d,"i'"]\\
X&X'\arrow[l,"f"]
\end{tikzcd}
\]
有如下\Term{canonical-isomorphism}{典范同构}，它们是\Term{duality}{对偶性}（SGA 4 XVIII 3.1.10）的平凡特例：
\[
\mathrm R\mathscr Hom(F,G)\simeq\mathrm Ri_*i^*G,\qquad
\mathrm R\mathscr Hom(f^*F,f^*G)\simeq\mathrm Ri'_*g^*i^*G.
\]
经这些\Term{isomorphism}{同构}，\Term{canonical}{典范}\Term{arrow}{箭头}
\[
f^*\mathrm R\mathscr Hom(F,G)\longrightarrow\mathrm R\mathscr Hom(f^*F,f^*G)
\]
等同于\Term{base-change}{基变换}\Term{arrow}{箭头}（SGA 4 XVII，原书未填细目）
\[f^*\mathrm Ri_*i^*G\longrightarrow\mathrm Ri'_*g^*i^*G.\]
\SourcePage{51}
而通过取\Term{limit}{极限}（SGA 4 VII 5.11），立即看出后者为\Term{isomorphism}{同构}。
\end{proof}
\begin{SGAresult}{\Term{lemma}{引理}}{4.4}
设 $X$ 为\Term{prescheme}{预概形}，$K\in\operatorname{ob}D_c^+(X)$ 满足 (1.7) 的条件 (i)、(ii)。$K$ 为\Term{dualizing-complex}{对偶化复形}，当且仅当对每个\Term{strict-localization}{严格局部化} $f:\overline X(\bar x)\to X$，$f^*K$ 为\Term{dualizing-complex}{对偶化复形}。
\end{SGAresult}
\begin{proof}
利用\Term{strict-localization}{严格局部化} $\overline X(\bar x)$ 上的每个\Term{constructible}{可构造}\Term{sheaf}{层}，都由 $X$ 上某个带 $\bar x$ 标点的\Term{etale}{平展}\Term{scheme}{概形}上的\Term{constructible}{可构造}\Term{sheaf}{层}诱导（SGA 4 IX 2.7.4），并应用 (4.3)。
\end{proof}
\textbf{4.5.} 设 $X$ 为\Term{prescheme}{预概形}，$\bar x$ 为 $X$ 的\Term{geometric-point}{几何点}，$f:X'=\overline X(\bar x)\to X$ 为相应的\Term{strict-localization}{严格局部化}。设 $Y$ 为 $x$ 在 $X$ 中的\Term{closure}{闭包}所给出的\Term{integral}{整}闭子\Term{prescheme}{预概形}；记 $g:\bar x\to Y$ 和 $i:Y\to X$ 为\Term{canonical}{典范}\Term{arrow}{箭头}，作\Term{cartesian-square}{笛卡尔方块}
\[
\begin{tikzcd}
Y\arrow[r,"i"]&X\\
\bar x\arrow[u,"g"]\arrow[r,"i'"']&X'\arrow[u,"f"']
\end{tikzcd}
\]
（$\bar x$ 同时是 $X'$ 的\Term{closed-point}{闭点}及 $Y$ 在 $\bar x$ 处的\Term{strict-localization}{严格局部化}。）记 $\mathrm R\Gamma_{\bar x}:D^+(X)\to D^+(\bar x)$ 为下式定义的\Term{functor}{函子}：
\[
\mathrm R\Gamma_{\bar x}=g^*\mathrm Ri^!.\tag{4.5.1}
\]
\SourcePage{52}
利用\Term{functor}{函子} $\mathrm Ri^!$ 和 $\mathrm Rj_*j^*$ 间的熟知关系，其中 $j:U\to X$ 为 $Y$ 的互补\Term{open-subset}{开集}，并取\Term{limit}{极限}（SGA 4 VII 5.11），得到\Term{canonical-isomorphism}{典范同构}
\[
\mathrm R\Gamma_{\bar x}\simeq\mathrm Ri'^!f^*.\tag{4.5.2}
\]
当 $x$ 为\Term{residue-field}{剩余域}\Term{separably-closed}{可分闭}的\Term{closed-point}{闭点}时，$\mathrm R\Gamma_{\bar x}$ 与通常的\Term{functor}{函子} $\mathrm R\Gamma_x$（SGA 4 V 4.3）一致，这也说明了这一记号的合理性。
设 $K_X\in\operatorname{ob}D^+(X)$，置
\[
K_Y=\mathrm Ri^!K_X,\qquad K_{X'}=f^*K_X,\qquad K_{\bar x}=\mathrm R\Gamma_{\bar x}(K_X).
\]
（注意，$K_{\bar x}$ 并不是 $K_X$ 在 $\bar x$ 处的\Term{fiber}{纤维}！）假定 $K_X$ 为\Term{dualizing-complex}{对偶化复形}。由 (1.15) 和 (4.4)，$K_Y$、$K_{X'}$、$K_{\bar x}$ 亦然。所以（2.1）$K_{\bar x}\simeq\Lambda[d]$，其中 $d\in\mathbb Z$。若 $d=0$ 且选择了\Term{isomorphism}{同构} $K_{\bar x}\simeq\Lambda$，就称 $K_X$ 在 $\bar x$ 处\Term{normalized}{正规化}。当 $x$ 为\Term{residue-field}{剩余域}\Term{separably-closed}{可分闭}的\Term{closed-point}{闭点}时，便得到 (4.1) 引进的概念。
设 $K_X$ 为\Term{dualizing-complex}{对偶化复形}，置 $D_X=\mathrm R\mathscr Hom(-,K_X)$、$D_Y=\mathrm R\mathscr Hom(-,K_Y)$ 等。设 $F\in\operatorname{ob}D_c^b(X)$，计算 $(D_XF)_{\bar x}$。有
\[
\begin{aligned}
(D_XF)_{\bar x}&\simeq g^*i^*D_XF\\
&\simeq g^*D_Y\mathrm Ri^!F&&\text{（由互补诱导公式 (4.2.1)）}\\
&\simeq D_{\bar x}g^*\mathrm Ri^!F&&\text{（由 (4.3)）},
\end{aligned}
\]
即
\[
(D_XF)_{\bar x}\simeq D_{\bar x}\mathrm R\Gamma_{\bar x}F.\tag{4.5.3}
\]
当然，也可经由
\SourcePage{53}
$X'$ 而非 $Y$ 作类似计算，得到\Term{canonical-isomorphism}{典范同构}
\[(D_XF)_{\bar x}\simeq D_{\bar x}\mathrm Ri'^!f^*F,\]
通过 (4.5.2) 的等同，它与 (4.5.3) 相容。
若 $K_X$ 还在 $\bar x$ 处\Term{normalized}{正规化}，则得到推广 (4.2.2)$'$ 和 (4.2.2)$''$ 的\Term{perfect-pairing}{完美配对}：
\[
D_X(F)_{\bar x}\times\mathrm R\Gamma_{\bar x}(F)\longrightarrow\Lambda,\tag{4.5.3$'$}
\]
\[
\mathscr H^i(D_XF)_{\bar x}\times H_{\bar x}^{-i}(F)\longrightarrow\Lambda,\tag{4.5.3$''$}
\]
其中置 $H_{\bar x}^i=\mathscr H^i\mathrm R\Gamma_{\bar x}$。
反之，设 $K_X\in\operatorname{ob}D^+(X)$ 满足 (1.7) 的条件 (i)、(ii)，且对 $X$ 的每个\Term{geometric-point}{几何点} $\bar x$，$K_{\bar x}=\mathrm R\Gamma_{\bar x}(K_X)$ 为\Term{dualizing-complex}{对偶化复形}。则对每个 $\bar x$，可以定义（参见 1.12 b) (ii)）\Term{canonical}{典范}\Term{arrow}{箭头}
\[
D_X(F)_{\bar x}\longrightarrow D_{\bar x}\mathrm R\Gamma_{\bar x}(F),\qquad F\in\operatorname{ob}D_c^b(X).\tag{4.5.3}
\]
若它对每个 $\bar x$ 和每个 $F$ 均为\Term{isomorphism}{同构}，则 $K_X$ 为\Term{dualizing-complex}{对偶化复形}。事实上，须验证\Term{canonical}{典范}\Term{arrow}{箭头} $F\to D_XD_XF$ 对每个 $F\in\operatorname{ob}D_c^b(X)$ 为\Term{isomorphism}{同构}。在任意\Term{geometric-point}{几何点} $\bar x$ 处，有
\[
\begin{aligned}
F_{\bar x}&\simeq D_{\bar x}D_{\bar x}(F_{\bar x})&&\text{（因为 $K_{\bar x}$ 为对偶化的）}\\
&\simeq D_{\bar x}\mathrm R\Gamma_{\bar x}D_X(F)&&\text{（由 (4.3) 及诱导公式 1.12 b) (i)）}\\
&\simeq D_XD_X(F)_{\bar x}&&\text{（由假设）}.
\end{aligned}
\]
\SourcePage{54}
因此，只需核对相容性便有 $(F,K_X)$ 为\Term{bidualizing}{双对偶化}的，从而证明断言。归纳起来：
\begin{SGAresult}{命题}{4.5.4}
设 $X$ 为\Term{prescheme}{预概形}，$K_X\in\operatorname{ob}D^+(X)$ 满足 (1.7) 的条件 (i)、(ii)。则 $K_X$ 为\Term{dualizing-complex}{对偶化复形}，当且仅当对 $X$ 的每个\Term{geometric-point}{几何点} $\bar x$，\Term{complex}{复形} $K_{\bar x}$ 为\Term{dualizing-complex}{对偶化复形}，且\Term{canonical}{典范}\Term{arrow}{箭头} (4.5.3) 对每个 $F\in\operatorname{ob}D_c^b(X)$ 都是\Term{isomorphism}{同构}。
\end{SGAresult}
\textbf{注记 4.5.5。} 在某种意义上，\Term{dualizing-complex}{对偶化复形}的存在应被视为整体\Term{duality-theorem}{对偶定理}的一项重要补充。后者说，对于一个\Term{separated}{分离}\Term{finite-type}{有限型}\Term{morphism}{态射} $f:X\to Y$（$Y$ \Term{quasi-compact}{拟紧}），存在\Term{canonical-isomorphism}{典范同构}
\[D_Y\mathrm Rf_!(F)\xrightarrow\sim\mathrm Rf_*D_X(F),\]
记号同 1.12。利用这一点计算 $\mathrm Rf_!(F)$（$F\in\operatorname{ob}D_c^b(X)$）需要掌握 $D_X(F)$ 的信息。当 $K_X$ 为\Term{dualizing-complex}{对偶化复形}时，\Term{isomorphism}{同构} (4.5.3) 提供这些信息；至少在理论上，它允许把 $D_X(F)$ 的\Term{fiber}{纤维}作为“纯\Term{space}{空间}”的\Term{cohomological-invariant}{上同调不变量}来计算。
\textbf{习题 4.5.6。} 设 $X$ 为带有\Term{dualizing-complex}{对偶化复形}（1.7）的\Term{prescheme}{预概形}。\Term{functor}{函子}族（参见 (4.5.3)$''$）
\[H_{\bar x}^i:D_c^b(X)\longrightarrow(\Lambda\text{-模}),\]
其中 $i$ 遍历 $\mathbb Z$、$\bar x$ 遍历 $X$ 的\Term{geometric-point}{几何点}，是
\SourcePage{55}
“\Term{conservative-functor}{保守}”的，即 $D_c^b(X)$ 中一个\Term{morphism}{态射} $F\to G$，若它诱导的 $H_{\bar x}^i(F)\to H_{\bar x}^i(G)$ 都是\Term{isomorphism}{同构}，则它本身是\Term{isomorphism}{同构}。
\textbf{4.6.} 设 $X$ 为\Term{prescheme}{预概形}，$j:Y\to X$ 为闭子\Term{prescheme}{预概形}，$i:\bar x\to Y$ 为 $Y$ 的\Term{geometric-point}{几何点}。又设 $K_X$ 为 $X$ 上在 $\bar x$ 处\Term{normalized}{正规化}的\Term{dualizing-complex}{对偶化复形}（参见 4.5）。在 (4.5) 中已看到，若 $Y$ 是 $x$ 的\Term{closure}{闭包}，则对 $F\in\operatorname{ob}D_c^b(X)$，有\Term{perfect-pairing}{完美配对}
\[\mathrm Rj^!(F)_{\bar x}\times D_X(F)_{\bar x}\longrightarrow\Lambda.\]
在一般情形中定义类似\Term{pairing}{配对}并不困难。引进 $i$ 的\Term{resolution}{分解}
\[\bar x\xrightarrow{i'}\overline{\{x\}}\xrightarrow{i''}Y.\]
对 $F\in\operatorname{ob}D_c^b(X)$，有\Term{canonical-isomorphism}{典范同构}
\[
\begin{aligned}
\mathrm R\Gamma_{\bar x}j^*D_X(F)
&\simeq\mathrm R\Gamma_{\bar x}D_Y\mathrm Rj^!(F)
&&\text{（互补诱导公式 (4.2.1)）}\\
&\simeq i'^*\mathrm Ri''^!D_Y\mathrm Rj^!(F)
&&\text{（$\mathrm R\Gamma_{\bar x}$ 的定义 (4.5.1)）}\\
&\simeq i'^*D_{\overline{\{x\}}}i''^*\mathrm Rj^!(F)
&&\text{（普通诱导 1.12 b) (i)）}\\
&\simeq D_{\bar x}i^*\mathrm Rj^!(F)&&\text{（由 (4.3)）},
\end{aligned}
\]
最终即\Term{canonical-isomorphism}{典范同构}
\[
\mathrm Rj^!(F)_{\bar x}\xrightarrow\sim D_{\bar x}\mathrm R\Gamma_{\bar x}j^*D_X(F).\tag{4.6.1}
\]
由此得到\Term{perfect-pairing}{完美配对}
\SourcePage{56}
\[
\mathrm Rj^!(F)_{\bar x}\times\mathrm R\Gamma_{\bar x}j^*D_X(F)\longrightarrow\Lambda,\tag{4.6.1$'$}
\]
\[
\mathscr H_Y^i(F)_{\bar x}\times H_{\bar x}^{-i}(D_X(F)|Y)\longrightarrow\Lambda.\tag{4.6.1$''$}
\]
\textbf{例 4.6.2。} 设 $X$ 为\Term{regular}{正则}\Term{prescheme}{预概形}，$\bar x$ 为 $X$ 的\Term{geometric-point}{几何点}，$d=\dim\mathscr O_{X,x}$。置 $K_X=(\mu_n)_X^{\otimes d}[2d]$。若 $X$ 满足 $(\mathrm{reg})$（3.1.1）及\Term{purity}{纯性}\Term{theorem}{定理}（3.1.4），则 $K_X$ 在 $\bar x$ 处\Term{canonical}{典范}地\Term{normalized}{正规化}。故若 $K_X$ 为\Term{dualizing-complex}{对偶化复形}，则对每个包含 $x$ 的\Term{closed-subset}{闭集} $Y$ 和每个 $F\in\operatorname{ob}D_c^b(X)$，有\Term{perfect-pairing}{完美配对}
\[
\mathscr H_Y^i(F)_{\bar x}\times H_{\bar x}^{2d-i}(F'|Y)\longrightarrow\Lambda,
\qquad F'=\mathrm R\mathscr Hom(F,(\mu_n)_X^{\otimes d}).
\]
最后给出 (4.6.1) 的一个\Term{strictly-local}{严格局部}变式。

\textbf{4.7.} 设 $X$ 为\Term{noetherian}{诺特}\Term{strictly-local}{严格局部}\Term{scheme}{概形}（SGA 4 VIII 4.2），\Term{closed-point}{闭点}为 $x$。设 $Y$ 为 $X$ 的\Term{nonempty}{非空}\Term{closed-subset}{闭集}，考虑\Term{canonical}{典范}\Term{immersion}{浸入}
\[U=X-Y\xrightarrow{i}V=X-\{x\}\xrightarrow{j}X.\]
再设 $K_X$ 为在 $x$ 处\Term{normalized}{正规化}的\Term{dualizing-complex}{对偶化复形}（4.1）；它给出 $U,V,Y,x$ 上“相应的”\Term{dualizing-complex}{对偶化复形}（由假设，$D_x=\operatorname{Hom}(-,\Lambda)$）：
\[K_V=j^*K_X,\qquad K_U=i^*K_V,\qquad K_Y=\mathrm R\underline\Gamma_Y(K_X).\]
设 $G\in\operatorname{ob}D_c^b(U)$，置 $G'=D_U(G)$。我们要定义一个自然\Term{pairing}{配对}
\[
\mathrm R\Gamma(U,G)\times\mathrm R\Gamma(V,i_!G')[-1]\longrightarrow\Lambda,\tag{4.7.1}
\]
即一个\Term{functorial}{函子性}\Term{homomorphism}{同态}
\SourcePage{57}
\[
\mathrm R\Gamma(U,G)\mathbin{\mathop\otimes\limits^{\mathrm L}_{\Lambda}}\mathrm R\Gamma(V,i_!G')[-1]\longrightarrow\Lambda.\tag{4.7.1$'$}
\]
（环 $\Lambda=\mathbb Z/\ell^\nu\mathbb Z$ 在 $\nu\geq2$ 时的\Term{cohomological-dimension}{上同调维数}是无穷的，因此只知道定义\Term{tensor-product}{张量积}
\[
\mathop\otimes^{\mathrm L}:D^-(\Lambda)\times D^-(\Lambda)\longrightarrow D^-(\Lambda)
\]
及
\[
\mathop\otimes^{\mathrm L}:D(\Lambda)\times D^b(\Lambda)_{\mathrm{torf}}\longrightarrow D(\Lambda).
\]\footnote{（1977 年增补）不正确：事实上，很容易把 $\mathop\otimes^{\mathrm L}:D^-(\Lambda)\times D^-(\Lambda)\to D^-(\Lambda)$ 延拓为 $\mathop\otimes^{\mathrm L}:D^-(\Lambda)\times D(\Lambda)\to D(\Lambda)$。}
因此将适当假设的选择留给读者，以保证 (4.7.1)$'$ 中 $\mathop\otimes^{\mathrm L}$ 的两个因子属于 $D^b(\Lambda)$，或者其中一个因子的 \Term{tor-dimension}{Tor 维数}\Term{finite}{有限}。）
先注意，由互补\Term{duality}{对偶}公式（1.12 a) (ii)）有\Term{canonical-isomorphism}{典范同构}
\[
i_!(G')\xrightarrow\sim D_V(\mathrm Ri_*G).\tag{4.7.2}
\]
另一方面，有普通\Term{duality-isomorphism}{对偶同构}
\[
\mathrm R\operatorname{Hom}(\Lambda_U,G)\xrightarrow\sim\mathrm R\operatorname{Hom}(\Lambda_V,\mathrm Ri_*G).\tag{4.7.3}
\]
由 $D_V$ 的定义（不依赖于 $K_V$ 为\Term{dualizing}{对偶化}这一事实），有\Term{canonical}{典范}\Term{arrow}{箭头}
\[
\mathrm Ri_*G\mathbin{\mathop\otimes\limits^{\mathrm L}_{\Lambda}}D_V(\mathrm Ri_*G)\longrightarrow K_V.\tag{4.7.4}
\]
由 (4.7.2)、(4.7.3)、(4.7.4) 得到（\Term{canonical}{典范}）\Term{arrow}{箭头}
\[
\mathrm R\Gamma(U,G)\mathbin{\mathop\otimes\limits^{\mathrm L}_{\Lambda}}\mathrm R\Gamma(V,i_!G')\longrightarrow\mathrm R\operatorname{Hom}(\Lambda_V,K_V).\tag{4.7.5}
\]
但 $K_V=j^!K_X$，且有\Term{duality-isomorphism}{对偶同构}（SGA 4 3.1.10）：
\[
\mathrm R\operatorname{Hom}(\Lambda_V,K_V)\xrightarrow\sim\mathrm R\operatorname{Hom}(\Lambda_{V,X},K_X).\tag{4.7.6}
\]
最后，\Term{exact-sequence}{正合列}
\SourcePage{58}
\[
0\longrightarrow\Lambda_{V,X}\longrightarrow\Lambda_X\longrightarrow\Lambda_x\longrightarrow0\tag{4.7.7}
\]
给出一个\Term{degree}{次数}为 $+1$ 的\Term{homomorphism}{同态}
\[
\mathrm R\operatorname{Hom}(\Lambda_{V,X},K_X)\longrightarrow\mathrm R\operatorname{Hom}(\Lambda_x,K_X)\simeq\Lambda.\tag{4.7.8}
\]
复合 (4.7.5)、(4.7.6)、(4.7.8)，便得到所宣称的\Term{arrow}{箭头} (4.7.1)$'$。
现在证明\Term{pairing}{配对} (4.7.1) 是“\Term{perfect}{完美}”的，即\Term{arrow}{箭头} (4.7.1)$'$ 把 $\mathrm R\Gamma(V,i_!G')[-1]$ 等同于 $D_x\mathrm R\Gamma(U,G)$。事实上，可设
\[G=(ji)^*F,\qquad F\in\operatorname{ob}D_c^b(X).\]
置 $F'=D_XF$，$F'_{U,X}=(ji)_!G'$，有\Term{exact-sequence}{正合列}
\[
0\longrightarrow F'_{U,X}\longrightarrow F'\longrightarrow F'|Y\longrightarrow0,\tag{4.7.9}
\]
它给出\Term{distinguished-triangle}{特殊三角形}
\[
\begin{tikzcd}
&\mathrm R\Gamma_x(F'_{U,X})\arrow[dr]\\
\mathrm R\Gamma_x(F'|Y)\arrow[ur,"{+1}"]&&\mathrm R\Gamma_x(F')\arrow[ll]
\end{tikzcd}\tag{b}
\]
由 $\mathrm R\Gamma_x$ 的定义，有
\[\mathrm R\Gamma_x(F'_{U,X})=\mathrm R\operatorname{Hom}(\Lambda_x,F'_{U,X}).\]
故由 (4.7.7) 及 $(F'_{U,X})_x=\mathrm R\Gamma(X,F'_{U,X})=0$，得到\Term{canonical-isomorphism}{典范同构}
\[
\mathrm R\operatorname{Hom}(\Lambda_{V,X},F'_{U,X})[-1]\xrightarrow\sim\mathrm R\Gamma_x(F'_{U,X}).
\]
\SourcePage{59}
但有
\[
\begin{aligned}
\mathrm R\operatorname{Hom}(\Lambda_{V,X},F'_{U,X})
&\simeq\mathrm R\operatorname{Hom}(j_!\Lambda_V,j_!i_!G')\\
&\simeq\mathrm R\operatorname{Hom}(\Lambda_V,i_!G')&&\text{由对偶性（SGA 4 3.1.10）},
\end{aligned}
\]
故最终有\Term{canonical-isomorphism}{典范同构}
\[
\mathrm R\Gamma_V(i_!G')[-1]\xrightarrow\sim\mathrm R\Gamma_x(F'_{U,X}).\tag{4.7.10}
\]
另一方面，有\Term{exact-sequence}{正合列}
\[
0\longrightarrow\Lambda_{U,X}\longrightarrow\Lambda_X\longrightarrow\Lambda_Y\longrightarrow0,\tag{4.7.11}
\]
它给出\Term{distinguished-triangle}{特殊三角形}
\[
\begin{tikzcd}
&\mathrm R\Gamma(U,F|U)\arrow[dl,"{+1}"']\\
\mathrm R\Gamma_Y(F)\arrow[rr]&&\mathrm R\Gamma_X(F)\arrow[ul]
\end{tikzcd}\tag{a}
\]
由 (4.7.1) 和 (4.7.10)，有自然\Term{pairing}{配对}
\[
\mathrm R\Gamma(U,F|U)\times\mathrm R\Gamma_x(F'_{U,X})\longrightarrow\Lambda.\tag{4.7.12}
\]
由 (4.6.1)，有\Term{perfect-pairing}{完美配对}
\[
\mathrm R\Gamma_Y(F)\times\mathrm R\Gamma_x(F'|Y)\longrightarrow\Lambda.\tag{4.7.13}
\]
最后，由普通\Term{induction-formula}{诱导公式}（1.12 b) (i)），有\Term{perfect-pairing}{完美配对}
\[
\mathrm R\Gamma(X,F)\times\mathrm R\Gamma_x(F')\longrightarrow\Lambda.\tag{4.7.14}
\]
有耐心在这些“\Term{canonical}{典范}\Term{arrow}{箭头}”的迷宫中辨清方向的读者，将验证\Term{pairing}{配对} (4.7.12)、(4.7.13)、(4.7.14) 与三角形 (a)、(b) 的\Term{arrow}{箭头}“相容”；从而\Term{pairing}{配对} (4.7.12) 和 (4.7.1) 是\Term{perfect}{完美}的。
\SourcePage{60}
三角形 (a)、(b) 的\Term{pairing}{配对}给出两条“互为\Term{transpose}{转置}”的\Term{exact-sequence}{正合列}：
\[
\cdots\longrightarrow H_Y^i(X,F)\longrightarrow H^i(X,F)\longrightarrow H^i(U,F)\longrightarrow\cdots
\tag{4.7.15}
\]
\[
\cdots\longleftarrow H_x^{-i}(Y,F'|Y)\longleftarrow H_x^{-i}(X,F')\longleftarrow H^{-i-1}(V,F'_{U,V})\longleftarrow\cdots.
\]
归纳起来，可陈述为：
\begin{SGAresult}{命题}{4.7.16}
设 $X$ 为\Term{noetherian}{诺特}\Term{strictly-local}{严格局部}\Term{scheme}{概形}，\Term{closed-point}{闭点}为 $x$。设 $Y$ 为 $X$ 的\Term{nonempty}{非空}\Term{closed-subset}{闭集}，$U=X-Y$，$V=X-x$，$i:U\to V$ 及 $j:V\to X$ 为\Term{canonical}{典范}包含。设 $K_X$ 为 $X$ 上在 $x$ 处\Term{normalized}{正规化}的\Term{dualizing-complex}{对偶化复形}（从而有 $V,U,Y$ 上相应的\Term{dualizing-complex}{对偶化复形}）。
对 $G\in\operatorname{ob}D_c^b(U)$，有\Term{perfect-pairing}{完美配对}
\[
\mathrm R\Gamma(U,G)\times\mathrm R\Gamma(V,i_!G')[-1]\longrightarrow\Lambda,\qquad G'=D_U(G).\tag{*}
\]
若 $G=F|U$，其中 $F\in\operatorname{ob}D_c^b(X)$，则有\Term{canonical-isomorphism}{典范同构}
\[
\mathrm R\Gamma(V,i_!G')[-1]\xrightarrow\sim\mathrm R\Gamma_x(F'_{U,X}),\tag{**}
\]
其中 $F'=D_XF$，$F'_{U,X}=(ji)_!(F'|U)=(ji)_!G'$，并有\Term{distinguished-triangle}{特殊三角形}
\[
\begin{tikzcd}[column sep=small]
&\mathrm R\Gamma(U,F|U)\arrow[dl,"{+1}"']&&&\mathrm R\Gamma_x(F'_{U,X})\arrow[dr]\\
\mathrm R\Gamma_Y(F)\arrow[rr]&&\mathrm R\Gamma(X,F)\arrow[ul]&\mathrm R\Gamma_x(F'|Y)\arrow[ur,"{+1}"]&&\mathrm R\Gamma_x(F')\arrow[ll]
\end{tikzcd}
\]
之间的\Term{perfect-pairing}{完美配对}，
\SourcePage{61}
与\Term{pairing}{配对} (*) 及\Term{isomorphism}{同构} (**) 相容，并给出两条互为\Term{transpose}{转置}的\Term{exact-sequence}{正合列} (4.7.15)。
\end{SGAresult}
\textbf{注记 4.7.17。} 在特殊情形 $Y=\{x\}$ 中，只需假定 $K_X$ 在 $x$ 处\Term{normalized}{正规化}（而不必为\Term{dualizing}{对偶化}的），即给定\Term{isomorphism}{同构} $\mathrm R\Gamma_x(K_X)\simeq\Lambda$，便可定义\Term{pairing}{配对} (4.7.1)：事实上，原来用到\Term{biduality}{双对偶性}的\Term{isomorphism}{同构} (4.7.2) 此时化为恒等式。由于\Term{pairing}{配对} (4.7.14) 总是\Term{perfect}{完美}的，可以看出 $X$ 上的\Term{local-duality}{局部对偶性}（以 (4.2.2)$'$ 的\Term{form}{形式}）\Term{equivalence}{等价}于 $U=X-\{x\}$ 上的“整体”\Term{duality}{对偶性}，即\Term{pairing}{配对} (4.7.1)
\[
\mathrm R\Gamma(U,G)\times\mathrm R\Gamma(U,G')[-1]\longrightarrow\Lambda
\]
为\Term{perfect}{完美}的，或\Term{equivalence}{等价}地，\Term{pairing}{配对}
\[H^i(U,G)\times H^{-i-1}(U,G')\longrightarrow\Lambda\]
为\Term{perfect}{完美}的。
假定在 $X$ 上有一个在 $x$ 处\Term{normalized}{正规化}的\Term{dualizing-complex}{对偶化复形} $K_X$。若 $X$ 为\Term{excellent}{优秀}的，且 $U$ 为 $d$ 维\Term{regular}{正则}\Term{prescheme}{预概形}（因此 $x$ 可以是一个\Term{isolated-singularity}{孤立奇点}），应能验证 $K_U=K_X|U$ \Term{canonical-isomorphism}{典范同构}于 $(\mu_n)_U^{\otimes d}[2d]$（无论如何，当 $X$ 为\Term{characteristic-zero}{特征零}域上\Term{locally-finite-type}{局部有限型}\Term{prescheme}{预概形}的\Term{strict-localization}{严格局部化}时，很容易验证）。于是 $X$ 上的\Term{local-duality}{局部对偶性}蕴含：对每个\Term{locally-constant}{局部常值}\Term{constructible}{可构造} $\Lambda_U$-模 $G$，有\Term{perfect-pairing}{完美配对}
\[
H^i(U,G)\times H^{2d-1-i}(U,G^\circ)\longrightarrow\Lambda,
\qquad G^\circ=\mathscr Hom(G,(\mu_n)_U^{\otimes d}),
\]
它在\Term{form}{形式}上对应于一个 $2d-1$ 维\Term{scheme}{概形}上的 \Term{poincare-duality}{Poincaré 对偶性}。
\SourcePage{62}
\SGAsection{5}{曲线上的局部对偶性}
\begin{SGAresult}{\Term{theorem}{定理}}{5.1}
设 $X$ 为\Term{dimension}{维数} $1$ 的\Term{noetherian}{诺特}\Term{regular}{正则}\Term{prescheme}{预概形}。假定 $n$ 与 $X$ 的各\Term{residue-characteristic}{剩余特征}互素。则\Term{purity}{纯性}\Term{theorem}{定理}在 $X$ 上成立（3.1.4），\Term{complex}{复形} $\Lambda_X$ 是\Term{dualizing-complex}{对偶化复形}（1.7）。
\end{SGAresult}
\begin{proof}
先假定已证明\Term{purity}{纯性}。$X$ 显然满足条件 $(\mathrm{reg})$（3.1.1），且由 SGA 4 X 2.2，条件 $(\mathrm{cdloc})_n$ 成立，所以由 (3.2.1) 得到 $\operatorname{dim.q.inj}(\Lambda_X)=2$。另一方面，由 (3.3.1 b) (i))（它的证明以\Term{purity}{纯性}为前提），(1.7) 的条件 (ii) 成立。因此由 (4.4)，只须在 $X$ 为\Term{strictly-local}{严格局部}时证明 $\Lambda_X$ 为\Term{dualizing-complex}{对偶化复形}。而为证明\Term{purity}{纯性}，由于它在\Term{closed-point}{闭点}\Term{neighborhood}{邻域}内是\Term{local-scheme}{局部}问题，同样可以假定 $X$ 为\Term{strictly-local}{严格局部}。最终，归结为在 $X$ 为\Term{strictly-local}{严格局部}时证明\Term{theorem}{定理}。
因此设 $X$ 为\Term{strictly-local}{严格局部}。设 $x$ 为其\Term{closed-point}{闭点}，$\pi=\operatorname{car}(k(x))$，$\eta$ 为其\Term{generic-point}{泛点}，$U=X-x=\operatorname{Spec}(k(\eta))$。设 $G=\operatorname{Gal}(\overline{k(\eta)}/k(\eta))$ 为 $U$ 的\Term{fundamental-group}{基本群}。已知（[CL] 第 IV 章第 2 节习题 2）有\Term{canonical}{典范}\Term{exact-sequence}{正合列}
\[
1\longrightarrow P\longrightarrow G\longrightarrow H\longrightarrow1,\tag{5.1.1}
\]
其中
\[
H=\prod_{\ell\ne\pi}\mathbb Z_\ell(1),\qquad
\mathbb Z_\ell(1)=\varprojlim_\nu\mu_{\ell^\nu},
\]
而 $P$ 为一个 pro-$\pi$ 群。（当然，$H$ \Term{isomorphism}{同构}于 $\prod_{\ell\ne\pi}\mathbb Z_\ell$，但不是\Term{canonical-isomorphism}{典范同构}。）
另一方面，已知（SGA 4 VIII 2）$U$ 的\Term{etale}{平展}\Term{cohomology}{上同调}可以解释为 $G$ 的 Galois \Term{cohomology}{上同调}；$\Lambda_U$-模\Term{sheaf}{层}（分别地，\Term{constructible}{可构造}\Term{sheaf}{层}）
\SourcePage{63}
对应于带有 $G$ \Term{continuous-action}{连续作用}的 $\Lambda$-模（分别地，\Term{finite-type}{有限型} $\Lambda$-模）。
于是，$\mathscr H_x^i(\Lambda_X)$ 由\Term{exact-sequence}{正合列}
\[
0\longrightarrow H_x^0(\Lambda_X)\longrightarrow H^0(X,\Lambda_X)\longrightarrow H^0(U,\Lambda_U)\longrightarrow H_x^1(\Lambda_X)\longrightarrow0,\tag{5.1.2}
\]
及\Term{isomorphism}{同构}
\[
H_x^i(\Lambda_X)\simeq H^{i-1}(U,\Lambda_U)\quad(i\geq2)\tag{5.1.3}
\]
确定。显然 $H^0(X,\Lambda_X)=H^0(U,\Lambda_U)=\Lambda$，因此 $H_x^0(\Lambda_X)=H_x^1(\Lambda_X)=0$。
另一方面，对一个（Galois）$G$-$\Lambda$-模 $M$，有 Hochschild—Serre \Term{spectral-sequence}{谱序列}
\[
E_2^{pq}=H^p(H,H^q(P,M))\Longrightarrow H^*(G,M).
\]
由于 $n$ 与 $\pi$ 互素而 $P$ 为 pro-$\pi$ 群，有 $H^q(P,M)=0$（$q\ne0$），故 $q\ne0$ 时 $E_2^{pq}=0$，从而有\Term{canonical-isomorphism}{典范同构}
\[
H^p(G,M)\simeq H^p(H,M^P).\tag{5.1.4}
\]
此外，由类似论证可知，当因子 $\mathbb Z_\pi(1)$ 在 $N$ 上\Term{trivial-action}{平凡作用}时，$H^p(H,N)$ 等同于 $H^p(\widehat{\mathbb Z}(1),N)$。因此，在上述约定下最终有\Term{canonical-isomorphism}{典范同构}
\[
H^p(G,M)\simeq H^p(\widehat{\mathbb Z}(1),M^P).\tag{5.1.5}
\]
而 $\widehat{\mathbb Z}$ 的 Galois \Term{cohomology}{上同调}是熟知的（[CG] 第 I 章，第 31 页）：对 $G$-$\Lambda$-模 $N$，有
\[
\begin{aligned}
H^i(\widehat{\mathbb Z},N)&=0&& (i\geq2),\\
H^0(\widehat{\mathbb Z},N)&=N^{\widehat{\mathbb Z}},\\
H^1(\widehat{\mathbb Z},N)&=N_{\widehat{\mathbb Z}}.
\end{aligned}\tag{5.1.6}
\]
\SourcePage{64}
由 (5.1.3)、(5.1.5)、(5.1.6) 立即得到
\[
H^i(U,\Lambda_U)=0\quad(i\geq2),\qquad H^1(U,\Lambda_U)\simeq\Lambda.
\]
后一个\Term{isomorphism}{同构}不是\Term{canonical}{典范}的，因为它依赖于选择 $\widehat{\mathbb Z}(1)$ 与 $\widehat{\mathbb Z}$ 的等同。但利用 \Term{kummer-exact-sequence}{Kummer 正合列}，可以得到（SGA 4 XIX 1.3）一个\Term{canonical-isomorphism}{典范同构}（由“\Term{local-fundamental-class}{局部基本类}”给出）
\[
H^1(U,\mu_n)\simeq\Lambda.\tag{5.1.7}
\]
\Term{purity}{纯性}因此得证。
现在证明 $\Lambda_X$ 为\Term{dualizing-complex}{对偶化复形}。由 (4.7.17)，只须证明对每个\Term{locally-constant}{局部常值}\Term{constructible}{可构造} $\Lambda_U$-模 $M$，\Term{pairing}{配对}
\[
H^i(U,M)\times H^{1-i}(U,\mathscr Hom(M,\Lambda_U))\longrightarrow H^1(U,\Lambda_U)\simeq(\mu_n)_x^{\otimes(-1)}\tag{5.1.8}
\]
为\Term{perfect-pairing}{完美对偶}。我们将用 Galois \Term{cohomology}{上同调}解释这一\Term{pairing}{配对}。为此使用如下初等\Term{lemma}{引理}。
\end{proof}
\begin{SGAresult}{\Term{lemma}{引理}}{5.1.9}
以 $D=\operatorname{Hom}(-,\Lambda)$ 表示\Term{finite-type}{有限型} $\Lambda$-\Term{module-category}{模范畴}上的 Pontryagin \Term{duality}{对偶}\Term{functor}{函子}。设 $M$ 为\Term{finite-type}{有限型} $\Lambda$-模，群 $G$ 在其上 $\Lambda$-\Term{linear}{线性}地\Term{action}{作用}。
(i) 有\Term{canonical-isomorphism}{典范同构} $D(M^G)\simeq(DM)_G$。
(ii) 若 $G$ 为在 $M$ 上\Term{continuous-action}{连续作用}的\Term{profinite}{预有限}群，且 $G$ 的 pro-阶与 $n$ 互素，则\Term{canonical}{典范}\Term{arrow}{箭头} $M^G\to M_G$ 为\Term{isomorphism}{同构}。
\end{SGAresult}
\SourcePage{65}
\begin{proof}
由定义，
\[M^G=\varprojlim_{s\in G}\operatorname{Ker}(M\xrightarrow{s-1}M),\]
故
\[
D(M^G)\simeq\varinjlim_{s\in G}\operatorname{Coker}(DM\xrightarrow{s-1}DM)\simeq(DM)_G,
\]
这证明了 (i)。
证明 (ii)。设 $M'$ 为由\Term{exact-sequence}{正合列} $0\to M'\to M\to M_G\to0$ 定义的 $\Lambda$-模。因为 $G$ 的 pro-阶与 $n$ 互素，有 $H^1(G,M')=0$。\Term{cohomology}{上同调}\Term{exact-sequence}{正合列}因此蕴含\Term{arrow}{箭头} $M^G\to M_G$ 是\Term{epimorphism}{满态射}。施以\Term{functor}{函子} $D$ 并考虑 (i)，又\Term{pushout}{推出}它是\Term{monomorphism}{单态射}，证明完毕。
\end{proof}
\begin{proof}[\Term{theorem}{定理} 5.1 证明的继续]
对 $M$ 和 $P$ 应用 (5.1.9)，得到 $\operatorname{Hom}(M,\Lambda)^P$ \Term{canonical}{典范}等同于 $\operatorname{Hom}(M^P,\Lambda)$。若选定 $\widehat{\mathbb Z}(1)$ 到 $\widehat{\mathbb Z}$ 的\Term{isomorphism}{同构}，考虑到 (5.1.5)，\Term{pairing}{配对} (5.1.8) 写成
\[
H^i(\widehat{\mathbb Z},M^P)\times H^{1-i}(\widehat{\mathbb Z},D(M^P))\longrightarrow H^1(\widehat{\mathbb Z},\Lambda)\simeq\Lambda.\tag{5.1.8$'$}
\]
由 (5.1.6) 和 (5.1.9 (i)) 立即得到 (5.1.8)$'$ 为\Term{perfect-pairing}{完美配对}。\Term{theorem}{定理} (5.1) 至此证毕。
\end{proof}

\SourcePage{66}
\textbf{参考文献}
\begin{description}
\item[{[2]}] Hironaka, H.：\emph{\Term{characteristic-zero}{特征零}域上\Term{algebraic-variety}{代数簇}的\Term{resolution-singularities}{奇点消解}}（Resolution of singularities of an algebraic variety over a field of characteristic zero），Ann. of Math.，第 79 卷（1964），第 109 页。
\item[{[1]}] Abhyankar, S.：\emph{\Term{characteristic}{特征} $p\ne0$ 的\Term{base-field}{基域}上\Term{algebra}{代数}\Term{surface}{曲面}的\Term{local-scheme}{局部}一致化}（Local uniformization on algebraic surfaces over ground fields of characteristic $p\ne0$），Ann. of Math.，第 63 卷（1956），第 491—526 页；以及\emph{算术\Term{surface}{曲面}的\Term{resolution-singularities}{奇点消解}}（Resolution of singularities of arithmetical surfaces），载 \emph{Arithmetical Algebraic Geometry}，Harper \& Row，纽约（1965）。
\item[{[H]}] Hartshorne, R.：\emph{\Term{residue}{留数}与\Term{duality}{对偶性}}（Residues and duality），Lecture Notes 第 20 号，Springer（1966）。
\item[{[SGA 4]}] IHÉS \Term{algebraic-geometry}{代数几何}讨论班：\emph{\Term{scheme}{概形}的\Term{etale-cohomology}{平展上同调}}（Cohomologie étale des schémas），M. Artin 和 A. Grothendieck，1963—1964。
\item[{[CL]}] Serre, J.-P.：\emph{\Term{local-scheme}{局部}域}（Corps locaux），Hermann（1962）。
\item[{[CG]}] Serre, J.-P.：\emph{Galois \Term{cohomology}{上同调}}（Cohomologie galoisienne），Lecture Notes 第 5 号，Springer（1964）。
\end{description}
